\documentclass[11pt]{article}

\usepackage[margin=1in]{geometry}
\usepackage{amsmath,amssymb,amsthm,mathtools}
\usepackage{enumitem}
\usepackage{graphicx}
\usepackage{microtype}
\usepackage[hidelinks]{hyperref}

\newtheorem{theorem}{Theorem}
\newtheorem{proposition}[theorem]{Proposition}
\newtheorem{lemma}[theorem]{Lemma}
\newtheorem{corollary}[theorem]{Corollary}
\theoremstyle{definition}
\newtheorem{definition}[theorem]{Definition}
\theoremstyle{remark}
\newtheorem{remark}[theorem]{Remark}

\newcommand{\E}{\mathcal E}
\newcommand{\Q}{\mathcal Q}
\newcommand{\VMO}{\mathrm{VMO}}
\newcommand{\ind}{\operatorname{ind}}
\newcommand{\Ind}{\operatorname{Ind}}
\newcommand{\wind}{\operatorname{wind}}
\newcommand{\Sym}{\operatorname{Sym}}
\newcommand{\tr}{\operatorname{tr}}
\newcommand{\dd}{\,\mathrm d}
\newcommand{\bd}{\partial}

\title{A VMO Morse Formula and a Complete Boundary-Extension\\
Criterion for Unoriented Line Fields on Surfaces}
\author{Candidate solution packet for arXiv:1407.1707}
\date{11 August 2026}

\begin{document}
\maketitle

\begin{center}
\fbox{\parbox{0.91\textwidth}{\textbf{Status.}
Candidate full solution, likely valid, awaiting expert review.  The theorem
settles the proposed extension in the compact oriented surface/Q-tensor
setting in which the source raises it.}}
\end{center}

\begin{abstract}
Canevari, Segatti, and Veneroni prove a Morse index formula for VMO vector
fields on manifolds with boundary and propose extending it to VMO line fields.
For a compact oriented surface $N$, the Q-tensors representing unoriented
lines form the oriented plane bundle
$\E=\Sym^2_0(TN)$.  Its transition functions have twice the angle of those of
$TN$, so $e(\E)=2e(TN)$.  We use this observation to define an intrinsic
half-integer line index and a boundary winding index for VMO data.  We prove
the exact formula
\[
  \ind_{\rm line}(Q,N)+\ind^{\rm line}_{\bd}(g)=\chi(N),
\]
and show that a bounded, uniformly positive VMO boundary line field $g$ has a
bounded, uniformly positive VMO extension if and only if
$\ind^{\rm line}_{\bd}(g)=\chi(N)$.  Thus the same invariant gives both the
Morse formula and a complete prescribed-boundary characterization.
\end{abstract}

\section{Source question and precise scope}

The source is G.~Canevari, A.~Segatti, and M.~Veneroni,
\emph{Morse's index formula in VMO for compact manifolds with boundary},
arXiv:1407.1707 \cite{CSV}.  At the start of its Q-tensor application in
Section~4 (source PDF page~17), after saying that it restricts to surfaces without boundary, it
states:
\begin{quote}
``Clearly, it would be interesting to extend Theorem [the VMO Morse theorem]
to VMO line fields.''
\end{quote}
The complete source passage is reproduced in Figure~\ref{fig:source}.

\begin{figure}[ht]
  \centering
  \IfFileExists{figures/open_problem_crop.png}{%
    \includegraphics[width=0.96\textwidth]{figures/open_problem_crop.png}%
  }{\fbox{\parbox{0.9\textwidth}{Source crop inserted after rendering the
  official arXiv source.}}}
  \caption{The proposal in Section~4 of arXiv:1407.1707.}
  \label{fig:source}
\end{figure}

The surrounding section works with a compact connected oriented surface
$N\subset\mathbb R^3$ and identifies a line field with a uniformly nonzero
section of its Q-tensor bundle.  We solve exactly the corresponding problem
when $N$ has boundary.  We also include the closed-surface specialization.
For higher-dimensional manifolds, a line field is a section of
$\mathbb P(TN)$ with fibre $\mathbb{RP}^{n-1}$; several obstruction groups can
occur, so no all-dimensional one-integer classification is claimed here.

The two explicit ``Question'' environments earlier in the source are not open
targets: Propositions in the same paper answer them.  The present packet
addresses only the separate line-field proposal quoted above.

\section{Definitions and the boundary invariant}

Let $N$ be a compact connected smooth oriented Riemannian surface.  Write
$P_x$ for the orthogonal projection onto $T_xN$.  Intrinsically, set
\[
  \E:=\Sym^2_0(TN),
\]
the bundle of trace-free symmetric endomorphisms of $TN$.  If
$N\subset\mathbb R^3$ is oriented with unit normal $\gamma$, this is precisely
the source's bundle
\[
 \Q_x=\{Q\in\Sym_0(3):Q\gamma(x)=0\}.
\]
For a unit tangent vector $n$, define
\begin{equation}\label{eq:qline}
 q(n):=\sqrt2\left(n\otimes n-\frac12P\right).
\end{equation}
Then $|q(n)|=1$ and $q(n)=q(-n)$.  Conversely, a unit element of $\E_x$
has a unique positive eigenspace, so the unit-circle bundle $S(\E)$ is
canonically the bundle of unoriented tangent lines.

Assume for now that $\bd N\ne\varnothing$.  Let $\nu$ be the outward unit
normal to $\bd N$ inside $TN$, and orient each boundary component by the unit
tangent $\tau$ for which $(\nu,\tau)$ is a positive frame of $TN$.  The
canonical boundary line section is
\begin{equation}\label{eq:qtau}
 q_\tau:=q(\tau)\in S(\E|_{\bd N}).
\end{equation}

If $g$ is a continuous nonzero section of $\E|_{\bd N}$, normalize it to
$\widehat g=g/|g|$.  The orientation and metric of $\E$ give a fibrewise
quarter-turn $J_\E$.  There is a unique map $r_g:\bd N\to S^1$ such that
\begin{equation}\label{eq:relativephase}
 \widehat g=(\Re r_g)q_\tau+(\Im r_g)J_\E q_\tau.
\end{equation}
Define the total relative winding, with the induced orientations on all
boundary components, by
\begin{equation}\label{eq:wcontinuous}
 w(g):=\sum_{C\subset\bd N}\deg(r_g|_C)\in\mathbb Z.
\end{equation}

For VMO data, we use the same averaging-and-fibre-projection convention as
the source.  Namely, embed $\E$ smoothly in a trivial Euclidean bundle, take
local averages of $g$ at scale $\varepsilon$, and project the result
orthogonally to $\E_x$.  Denote the resulting continuous section by
$g_\varepsilon$.  If
\begin{equation}\label{eq:bounds-g}
 g\in\VMO(\bd N,\E),\qquad 0<c\le |g|\le C<\infty\quad\text{a.e.},
\end{equation}
then $g_\varepsilon$ is nonzero for all sufficiently small $\varepsilon$, and
$w(g_\varepsilon)$ is independent of that small scale.  We therefore set
\begin{equation}\label{eq:wVMO}
 w(g):=w(g_\varepsilon).
\end{equation}
The bundle-VMO lemma below proves that this is well-defined.  We likewise
write $e(\E,g)$ for the stable relative Euler class
$e(\E,g_\varepsilon)$ when $g$ is VMO.

\begin{definition}[Line and boundary indices]\label{def:indices}
Let $Q\in\VMO(N,\E)$ have VMO trace $g$ satisfying
\eqref{eq:bounds-g}.  Let $Q_\varepsilon$ be a source-type continuous
regularization, nonzero on $\bd N$.  Its ordinary plane-bundle index is the
sum of the integer local degrees of a transverse perturbation, and is denoted
$\Ind_\E(Q,N)$.  Define
\begin{equation}\label{eq:lineindices}
 \ind_{\rm line}(Q,N):=\frac12\Ind_\E(Q,N),
 \qquad
 \ind^{\rm line}_{\bd}(g):=-\frac12w(g).
\end{equation}
Both may be half-integers.  This is the standard normalization in which a
half-strength nematic defect has line index $1/2$.
\end{definition}

\section{Idea of the proof}

The sign ambiguity that blocks a literal ``points inward'' definition is
removed by representing an unoriented line as the Q-tensor $q(n)$.  In an
oriented tangent frame, rotating $n$ by $\theta$ rotates $q(n)$ by
$2\theta$.  Thus the Q-bundle has twice the Euler class of $TN$.

Along the boundary, the tangent line gives a canonical reference section
$q_\tau$.  Relative to it, the boundary datum has an integer winding $w(g)$.
The relative Euler number of the Q-bundle is exactly
\[
  2\chi(N)+w(g).
\]
It is also the total ordinary index of any Q-section with boundary value $g$.
Dividing by two and moving the winding term to the left gives the desired
Morse formula.

The analytic VMO step introduces no new topological obstruction.  The
source's local-average projection works for any smooth Euclidean subbundle,
not only $TN$, so it applies to $\E$.  Uniform positivity prevents the
regularized Q-fields from acquiring zeros.  Conversely, if the relative Euler
number vanishes, average the VMO datum into a collar and use the vanishing
Euler obstruction to extend the resulting continuous nonzero section over the
two-dimensional core.  The glued field is VMO, has the prescribed trace, and
can be radially clamped to uniform positive upper and lower bounds.

\section{Topological calculation}

\begin{lemma}[The doubled Euler bundle]\label{lem:double}
The bundle $\E=\Sym^2_0(TN)$ is an oriented metric plane bundle.  With the
orientations above,
\begin{equation}\label{eq:relativeeulerqtau}
 \left\langle e(\E,q_\tau),[N,\bd N]\right\rangle=2\chi(N).
\end{equation}
For any continuous nonzero boundary section $g$,
\begin{equation}\label{eq:relativeeuler}
 \left\langle e(\E,g),[N,\bd N]\right\rangle
   =2\chi(N)+w(g).
\end{equation}
\end{lemma}

\begin{proof}
In a positively oriented orthonormal tangent frame $(e_1,e_2)$, use the
orthonormal frame
\begin{equation}\label{eq:Eframe}
 A_1=\frac{e_1\otimes e_1-e_2\otimes e_2}{\sqrt2},
 \qquad
 A_2=\frac{e_1\otimes e_2+e_2\otimes e_1}{\sqrt2}
\end{equation}
of $\E$.  If $(e_1,e_2)$ is rotated through $\theta$, direct expansion shows
that $(A_1,A_2)$ is rotated through $2\theta$.  Thus the oriented frame
transition maps of $\E$ are the squares of those of $TN$.

The relative Euler number of $(TN,\tau)$ is $\chi(N)$.  One way to see the
sign is to apply the continuous Morse formula from \cite{CSV} to a tangent
vector field whose boundary value is $\tau$: its inward boundary set is
empty, so its total zero index is $\chi(N)$.  Squaring every transition map
and the boundary clutching map doubles the obstruction cocycle.  Since
$q(\tau)$ is precisely the squared boundary section, this proves
\eqref{eq:relativeeulerqtau}.

Changing a nonzero boundary section from $q_\tau$ to $g$ changes the relative
Euler obstruction by the total degree of the transition map $r_g$ in
\eqref{eq:relativephase}.  With the induced boundary orientation, the sign is
positive.  For example, on the positively oriented disk $q_\tau$ has degree
$2$ in a constant trivialization, while a constant section has relative
winding $-2$ and relative Euler number $0$.  Hence
\[
 \langle e(\E,g),[N,\bd N]\rangle
 =\langle e(\E,q_\tau),[N,\bd N]\rangle+\deg(r_g),
\]
which is \eqref{eq:relativeeuler}.
\end{proof}

\begin{lemma}[Relative index and extension]\label{lem:obstruction}
Let $E\to M$ be an oriented metric plane bundle over a compact connected
oriented surface with boundary, and let $h$ be a nonzero continuous section
of $E|_{\bd M}$.
\begin{enumerate}[label=\textup{(\roman*)}]
\item The total local index of any continuous extension, after a transverse
perturbation fixed near the boundary, is
$\langle e(E,h),[M,\bd M]\rangle$.
\item The section $h$ extends to a continuous nowhere-zero section on $M$ if
and only if this integer is zero.
\end{enumerate}
\end{lemma}

\begin{proof}
Triangulate the pair $(M,\bd M)$ and normalize sections to the circle bundle
$S(E)$.  Since $S^1$ is path connected, $h$ extends over the relative
one-skeleton.  The winding around each oriented two-cell is the obstruction
to extending over that cell.  These windings form a relative $2$-cocycle with
values in $\pi_1(S^1)=\mathbb Z$; its class is the relative Euler class.
For a transverse section, the same cell degrees are the signed local zero
indices, proving (i).

There are no higher cells and no higher obstructions.  Thus a nonzero
extension exists exactly when the relative obstruction class vanishes.  By
Poincare--Lefschetz duality,
$H^2(M,\bd M;\mathbb Z)\cong\mathbb Z$ for connected oriented $M$, and its
evaluation on $[M,\bd M]$ detects the class.  This proves (ii).
\end{proof}

\section{The VMO passage}

We isolate the precise analytic input.  It is the source's construction with
the tangent bundle replaced by a smooth Euclidean subbundle.

\begin{lemma}[VMO bundle regularization and collar]\label{lem:vmo}
Let $E\to M$ be a smooth metric vector bundle embedded as a smooth subbundle
of $M\times\mathbb R^m$.
\begin{enumerate}[label=\textup{(\alph*)}]
\item A VMO section with VMO trace admits continuous local-average
regularizations obtained by orthogonal projection to $E_x$.  If the trace is
bounded away from zero, the regularizations are uniformly nonzero on the
boundary for small scales; their relative indices and boundary windings are
then independent of the scale.
\item If $s\in\VMO(M,E)$ satisfies $c\le |s|\le C$ almost everywhere, its
regularizations satisfy $c/2\le |s_\varepsilon|\le 2C$ everywhere for all
sufficiently small $\varepsilon$.
\item If $h\in\VMO(\bd M,E|_{\bd M})$ satisfies
$c\le |h|\le C$, then on a sufficiently thin collar
$\Phi:\bd M\times[0,r]\to M$ there is a VMO section $H$, continuous away
from $\bd M$, with trace $h$, with $H|_{\Phi(\bd M\times\{r\})}$ continuous,
and with $c/3\le |H|\le3C$.
\end{enumerate}
\end{lemma}

\begin{proof}
We indicate why the proofs in Sections~3--4 of \cite{CSV} apply without a
hidden tangent-vector assumption.  In local coordinates, average the ambient
$\mathbb R^m$-valued representative and apply the orthogonal projection
$P_{E_x}$.  The map $x\mapsto P_{E_x}$ is smooth.  The error made by replacing
$P_{E_y}$ by $P_{E_x}$ on a ball of radius $\varepsilon$ is therefore
$o(1)$ uniformly; the remaining error is exactly the VMO mean oscillation.
The trace construction on the double gives uniform compatibility between the
interior and boundary regularizations.  Homotopy invariance then makes the
indices constant for small scale.  These are precisely the ingredients used
for the tangent bundle in the source.

For (b), on every sufficiently small ball $B$,
\[
 |s_B|\ge \frac1{|B|}\int_B|s|-\frac1{|B|}\int_B|s-s_B|
 \ge c-o(1),
\]
uniformly in the centre, and smooth fibre projection changes this by another
$o(1)$.  The upper bound is analogous.

For (c), put $x=\Phi(y,t)$ and take
\[
 H(x)=P_{E_x}\left(\frac1{|B_t^{\bd M}(y)|}
                 \int_{B_t^{\bd M}(y)}h(z)\dd z\right),\qquad t>0.
\]
The VMO oscillation estimate gives the stated bounds and continuity for
$t>0$.  The standard VMO collar lemma proved in \cite{CSV} says that this
averaging extension has trace $h$ and is VMO across $t=0$; composition with
the smooth family $P_{E_x}$ preserves those properties.  This proves (c).
\end{proof}

\section{Morse formula and full characterization}

\begin{theorem}[VMO Morse formula for line fields]\label{thm:morse}
Let $N$ be a compact connected oriented smooth surface with nonempty boundary,
and put $\E=\Sym^2_0(TN)$.  Let $Q\in\VMO(N,\E)$ have VMO trace
$g\in\VMO(\bd N,\E)$ satisfying $0<c\le |g|\le C$.  Then the indices in
Definition~\ref{def:indices} are well-defined and satisfy
\begin{equation}\label{eq:morse}
 \boxed{\quad
 \ind_{\rm line}(Q,N)+\ind^{\rm line}_{\bd}(g)=\chi(N).
 \quad}
\end{equation}
Equivalently,
\begin{equation}\label{eq:integerformula}
 \Ind_\E(Q,N)=2\chi(N)+w(g).
\end{equation}
\end{theorem}

\begin{proof}
Lemma~\ref{lem:vmo}(a) gives a continuous regularization $Q_\varepsilon$ that
is nonzero on the boundary and has boundary value homotopic through nonzero
sections to $g_\varepsilon$.  By Lemma~\ref{lem:obstruction}(i), its total
integer index is
\[
 \Ind_\E(Q,N)=\langle e(\E,g_\varepsilon),[N,\bd N]\rangle.
\]
Lemma~\ref{lem:double} and the definition of the stable VMO winding give
\[
 \Ind_\E(Q,N)=2\chi(N)+w(g).
\]
Dividing by two and using \eqref{eq:lineindices} proves \eqref{eq:morse}.
\end{proof}

\begin{theorem}[Prescribed-boundary characterization]\label{thm:extension}
Under the hypotheses of Theorem~\ref{thm:morse}, assume only that the boundary
datum $g$ is given.  The following are equivalent:
\begin{enumerate}[label=\textup{(\roman*)}]
\item There exists $Q\in\VMO(N,\E)$ with trace $g$ and constants
$0<c'\le C'<\infty$ such that
\[
 c'\le |Q(x)|\le C'\quad\text{for almost every }x\in N.
\]
Thus $Q$ is a bounded, uniformly positive VMO line field.
\item The boundary line index satisfies
\begin{equation}\label{eq:criterion-index}
 \ind^{\rm line}_{\bd}(g)=\chi(N).
\end{equation}
\item The explicit winding condition
\begin{equation}\label{eq:criterion-winding}
 2\chi(N)+w(g)=0
\end{equation}
holds.
\end{enumerate}
\end{theorem}

\begin{proof}
The equivalence of (ii) and (iii) is the definition of the boundary index.
If (i) holds, Lemma~\ref{lem:vmo}(b) makes all sufficiently small
regularizations nonzero on the whole of $N$.  Hence
$\Ind_\E(Q,N)=0$, and Theorem~\ref{thm:morse} gives (ii).

Assume (iii).  Apply Lemma~\ref{lem:vmo}(c) to extend $g$ as a uniformly
nonzero VMO section $H$ over a collar $U_r$.  On the inner collar boundary
$C_r$, the section $h=H|_{C_r}$ is continuous and nonzero.  Let
$N_r=N\setminus U_r$.  Additivity of the relative Euler obstruction under
gluing and the absence of zeros in the collar give
\[
 \langle e(\E|_{N_r},h),[N_r,C_r]\rangle
 =\langle e(\E,g),[N,\bd N]\rangle
 =2\chi(N)+w(g)=0.
\]
By Lemma~\ref{lem:obstruction}(ii), $h$ extends over $N_r$ to a continuous
nowhere-zero section $K$.  Glue $K$ to $H$.  The result is VMO, has trace
$g$, and is uniformly positive and bounded because $N_r$ is compact and both
pieces are nonzero.

If desired, compose fibrewise with a Lipschitz radial clamping map that is the
identity on the essential range of $g$.  This preserves VMO and the trace and
gives any convenient finite positive bounds.  Thus (i) holds.
\end{proof}

\begin{corollary}[Closed surfaces]\label{cor:closed}
If $N$ is compact, connected, oriented, and has no boundary, every VMO
Q-field for which the line index is defined satisfies
\[
 \ind_{\rm line}(Q,N)=\chi(N).
\]
In particular, a bounded uniformly positive VMO line field exists if and only
if $\chi(N)=0$.
\end{corollary}

\begin{proof}
The ordinary total index of a section of $\E$ is
$\langle e(\E),[N]\rangle=2\chi(N)$ by the double-angle calculation in
Lemma~\ref{lem:double}.  Divide by two.  Necessity for a nonzero VMO line
field follows from Lemma~\ref{lem:vmo}(b), and sufficiency when $\chi(N)=0$
is the closed-surface rank-two obstruction statement (or the source's
continuous Poincare--Kneiser theorem, whose continuous section is VMO).
\end{proof}

\section{Checks, examples, and limitations}

\paragraph{Disk sign check.}
Trivialize $TN$ on the unit disk by the Cartesian frame.  Along the positively
oriented circle, $\tau(\theta)=(-\sin\theta,\cos\theta)$, so $q_\tau$ has
phase $2\theta+\pi$ and degree $2$.  A constant boundary Q-field has relative
winding $-2$ and extends constantly.  Formula \eqref{eq:relativeeuler} gives
$2\chi(D^2)-2=0$, as required.

\paragraph{Annulus check.}
The induced orientations on the two boundary circles are opposite, and
$\chi=0$.  The contributions of $q_\tau$ cancel.  A constant Q-field has
$w=0$ and extends, again agreeing with the criterion.

\paragraph{Gauge invariance.}
Changing an auxiliary tangent frame by angle $\alpha$ changes Q-coordinates
by $2\alpha$ for both $g$ and $q_\tau$.  Their relative phase, hence $w(g)$,
is unchanged.

\paragraph{Scope.}
The result is complete for the compact oriented surface/Q-tensor problem
posed by the surrounding source section.  It does not classify line fields on
nonorientable surfaces or arbitrary higher-dimensional manifolds.  The latter
is generally an obstruction-theory problem for $\mathbb P(TN)$ rather than a
single Euler-number formula.

\paragraph{Novelty check.}
Exact-phrase and close-variant web/arXiv searches, the run indexes, and the
complete OpenAlex citing-work list for the published source were checked on
11 August 2026.  No later explicit solution of this VMO boundary problem was
found.  The continuous relative-Euler ingredients are classical; the novelty
claim concerns their exact line-field normalization and the VMO
trace/extension theorem.  See \texttt{novelty\_search.md} for the search
bounds.

\paragraph{Computational sanity check.}
The included script \texttt{code/verify\_transition\_and\_winding.py} verifies
the double-angle matrix identity and the disk/annulus winding signs on a dense
mesh.  These computations are checks, not part of the proof.

\section{Human-review recommendation}

This packet should receive expert review as a full candidate.  The most useful
checks are:
\begin{enumerate}
\item verify the sign in
$\langle e(\E,g),[N,\bd N]\rangle=2\chi(N)+w(g)$ under the stated boundary
orientation;
\item confirm that the source's double/trace and collar arguments depend only
on a smooth Euclidean subbundle and therefore apply to $\E$;
\item verify the gluing identification of the relative Euler class on the
collar core;
\item compare the term ``extend Theorem to VMO line fields'' with the source
authors' intended defect-index normalization.
\end{enumerate}

\begin{thebibliography}{9}
\bibitem{CSV}
G.~Canevari, A.~Segatti, and M.~Veneroni,
\emph{Morse's index formula in VMO for compact manifolds with boundary},
Journal of Functional Analysis 269 (2015), 3043--3082;
arXiv:1407.1707.

\end{thebibliography}

\end{document}
